% Compile with pdfLaTeX from this directory or the repository root.
\makeatletter
\def\input@path{{../../../preamble/}{book2026-2027/preamble/}}
\makeatother
\input{preamble.tex}

\geometry{letterpaper,margin=0.72in}
\linespread{1.06}
\setlength{\parskip}{4pt minus 3pt}
\setcounter{secnumdepth}{1}
\renewcommand{\thesection}{Step \arabic{section}}
\setlist[itemize]{leftmargin=1.2em,topsep=3pt,itemsep=2pt,parsep=0pt}
\hypersetup{pdftitle={Confidence \& Interval Estimation | Grade 11 Term 2 Outline},
  linkcolor=taskLinkDarkBlue}

\newcolumntype{Y}{>{\raggedright\arraybackslash}X}
\newcommand{\VisualHeading}[1]{\par\smallskip\noindent\textbf{#1}\par\nobreak}
\newcommand{\Criterion}[2]{\par{\small\noindent\textbf{Criterion #1.}\enspace #2\par}}
\tikzset{
  outline interval/.style={draw=taskLinkDarkBlue,very thick},
  outline guide/.style={draw=black!65,dashed}
}

\begin{document}

\begin{center}
  {\LARGE\bfseries Confidence \& Interval Estimation\par}
  \smallskip
  {\large Grade 11 Global Economics \(\cdot\) Term 2\par}
\end{center}

\subsection*{What can a business infer about its average transaction processing cost and the share of transactions processed late?}

\textbf{Hypothetical case.} A business processed 10,000 online card transactions in October. A transaction counts as late if it missed the service deadline.

\section{From a sample to a population}

The population is the business's 10,000 October online card transactions. A point estimate gives one value, such as a hypothetical sample mean of \textbf{\$48}. An interval estimate gives a range of plausible values for the population mean, such as \textbf{\$46.04 to \$49.96} in Step 2.
\[
  \text{estimate} \pm \text{critical value} \times \text{standard error}
\]
The estimate sets the centre. The critical value reflects the confidence level. The standard error measures how much the estimate varies across samples. Their product (critical value \(\times\) standard error) is the \textbf{margin of error}, half the interval's width.

\VisualHeading{C1--C3 \(\cdot\) Population, sample and repeated-sampling coverage}

\noindent
\begin{minipage}[c]{0.50\linewidth}
\small
\setlength{\tabcolsep}{4pt}
\begin{tabularx}{\linewidth}{@{}YYY@{}}
  \toprule
  \textbf{Quantity} & \textbf{Population parameter} & \textbf{Sample statistic}\\
  \midrule
  Mean cost & \(\mu\) & \(\bar{x}\)\\
  Cost standard deviation & \(\sigma\) & \(s\)\\
  Late proportion & \(p\) & \(\hat{p}\)\\
  \bottomrule
\end{tabularx}
\end{minipage}\hfill
\begin{minipage}[c]{0.47\linewidth}
  \centering
  \small\textbf{Five illustrative confidence intervals from repeated samples}\par\smallskip
  % These intervals preserve the SVG's relative positions on an unlabeled scale.
  \begin{tikzpicture}[x=0.033cm,y=-0.033cm,font=\small]
    \node at (151,6) {Fixed \(\mu\)};
    \draw[outline guide] (151,16)--(151,145);
    \foreach \drawnum/\row/\low/\high in {1/34/100/220,2/58/119/239,3/82/92/212,5/136/110/230} {
      \node[anchor=east] at (80,\row) {Draw \drawnum};
      \draw[outline interval] (\low,\row)--(\high,\row);
    }
    \node[anchor=east] at (80,106) {Draw 4};
    \draw[outline interval,draw=darkred] (170,106)--(290,106);
    \node[below] at (230,106) {misses \(\mu\)};
  \end{tikzpicture}
\end{minipage}

\smallskip
Each bar represents a hypothetical draw in a separate illustration. The vertical line is the fixed population mean. Intervals 1, 2, 3 and 5 cross it. Interval 4 lies to its right and misses it. With repeated sampling under the same conditions, about 95\% of intervals made by a 95\% procedure would contain the fixed \(\mu\). Each calculated interval either contains \(\mu\) or misses it.

\Criterion{C1}{Distinguishes population parameters from sample statistics by interpreting and comparing the notation for population and sample means, standard deviations, and proportions in economic and financial contexts.}
\Criterion{C2}{Compares point and interval estimation and interprets confidence levels through their long-run coverage meaning, distinguishing this from the probability that a particular calculated interval contains the parameter.}
\Criterion{C3}{Identifies the components of a confidence interval, including point estimate, critical value, standard error, and margin of error, explaining how each component affects precision and reliability.}

\clearpage
\section{Build \texorpdfstring{\(z\)}{z} intervals for a mean and a proportion}

\textbf{Hypothetical sample A.} Draw a simple random sample of 100 of the 10,000 transactions without replacement. For the mean, assume \(\sigma=\$10\) was established independently for this population. Both 95\% intervals use \(z^*=1.96\).

\VisualHeading{C4--C5 \(\cdot\) Two 95\% \(z\) intervals from the same sample}

{\small\noindent
\begin{tabularx}{\linewidth}{@{}lYY@{}}
  \toprule
  \textbf{Quantity} & \textbf{Mean cost (C4)} & \textbf{Late share (C5)}\\
  \midrule
  Estimate & \(\bar{x}=\$48\) & \(\hat{p}=20/100=0.20\)\\[1pt]
  Standard error
    & \(\sigma/\sqrt{n}=\$10/\sqrt{100}=\$1\)
    & \(\sqrt{\hat{p}(1-\hat{p})/n}\)\newline
      \(=\sqrt{0.20\times 0.80/100}=0.0400\)\\[1pt]
  Margin of error
    & \(1.96\times \$1=\$1.96\)
    & \(1.96\times 0.04=0.0784\)\newline (7.84 percentage points)\\[1pt]
  95\% interval
    & \(\$48\pm \$1.96\)\newline \textbf{\$46.04 to \$49.96}
    & \(0.20\pm 0.0784\)\newline \textbf{12.16\% to 27.84\%}\\
  \bottomrule
\end{tabularx}\par}

\textbf{Check the conditions.} The sample is random and covers only 1\% of the population, supporting approximate independence. For the mean, \(\sigma\) is known and \(n=100\) supports a normal approximation assuming costs have no extreme skew or outliers. For the one-proportion \(z\) interval, 20 late and 80 on-time both meet the minimum of 10.

We are 95\% confident the October population mean cost lies between \textbf{\$46.04 and \$49.96}. For the late share, we are 95\% confident it lies between \textbf{12.16\% and 27.84\%}.

\Criterion{C4}{Constructs confidence intervals for a population mean with known variance using the normal distribution, applying them to contexts such as average income, average investment return, or average price level.}
\Criterion{C5}{Constructs confidence intervals for a population proportion using sample proportions, applying them to contexts such as poverty rates, defaulting borrowers, or successful investments.}

\section{Plan sample size and compare precision}

To plan a new sample, set the target margin of error \(\boldsymbol{E}\). For 95\% confidence, use \(z^*=1.96\).
\begin{itemize}
  \item \textbf{Mean, known \(\sigma\).} \(n\ge (z^*\sigma/E)^2\). With \(\sigma=\$10\) and \(E=\$2\), \((1.96\times 10/2)^2=96.04\), so sample \textbf{at least~97}.
  \item \textbf{Proportion.} \(n\ge (z^*/E)^2 p^*(1-p^*)\). With \(E=0.05\) (5 percentage points) and planning proportion \(p^*=0.50\), \((1.96/0.05)^2\times 0.50\times 0.50=384.16\), so sample \textbf{at least 385}. The choice 0.50 gives the largest planned sample because it maximises \(p^*(1-p^*)\).
\end{itemize}
Round up to a whole transaction so the planned margin of error stays at or below \(E\).

\VisualHeading{C6--C7 \(\cdot\) Compare widths with \(\bar{x}\) fixed at \$48 and \(\sigma\) fixed at \$10}

\noindent
\begin{minipage}[c]{0.605\linewidth}
  \centering
  \small\textbf{Three mean intervals on the same cost scale}\par\smallskip
  \begin{tikzpicture}[x=1.09cm,y=0.60cm,font=\small]
    % Horizontal coordinates are dollar differences from the common centre.
    \node[above] at (0,3.46) {Common centre \$48};
    \draw[outline guide] (0,0.55)--(0,3.46);
    \foreach \row/\conf/\size/\margin/\low/\high in {
      3/95/100/1.96/46.04/49.96,2/99/100/2.576/45.42/50.58,1/95/400/0.98/47.02/48.98
    } {
      \node[anchor=east] at (-3.89,\row) {\conf\% \(\cdot\) \(n=\size\)};
      \draw[outline interval] (-\margin,\row)--(\margin,\row);
      \draw[outline interval] (-\margin,\row-.15)--(-\margin,\row+.15)
        (\margin,\row-.15)--(\margin,\row+.15);
      \node[left=2pt] at (-\margin,\row) {\$\low};
      \node[right=2pt] at (\margin,\row) {\$\high};
    }
  \end{tikzpicture}
\end{minipage}\hfill
\begin{minipage}[c]{0.38\linewidth}
  \raggedright\small
  The first bar repeats sample A. Raising confidence to 99\% at \(n=100\) increases the critical value and widens the interval. A hypothetical \(n=400\) at 95\% halves the standard error and the width.
\end{minipage}

\Criterion{C6}{Determines the sample size required to estimate a population's mean or proportion with a specified margin of error, using it to plan economic surveys or financial studies with predetermined precision.}
\Criterion{C7}{Compares confidence intervals constructed from samples of different sizes and confidence levels, analysing the trade-off between confidence and precision.}

\clearpage
\section{Use Student's \texorpdfstring{\(t\)}{t} and check the design}

\textbf{Hypothetical sample B.} Take a separate simple random sample of 25 transactions from the same October population. In this example, assume \(\sigma\) is unknown. The sample gives \(\bar{x}=\$48\) and \(s=\$10\).

Apply the Term 1 distinction to intervals for a mean. With known \(\sigma\), use \(z\) as in Step 2. Here \(s\) estimates \(\sigma\), so use Student's \(t\) with \(\boldsymbol{df=25-1=24}\). For 95\% confidence, \(t^*\approx 2.064\) (2.5\% in each tail).

\noindent
\begin{minipage}[c]{0.39\linewidth}
  \(\mathrm{SE}=s/\sqrt{n}=\$2\)\par
  Margin \(=2.064\times \$2\approx \$4.13\)\par
  95\% interval for \(\mu\approx \$48\pm \$4.13\)\par
  \(=\$43.87\) to \(\$52.13\)
\end{minipage}\hfill
\begin{minipage}[c]{0.57\linewidth}
  \centering
  \small\textbf{C8--C9 \(\cdot\) Check the shape of the 25 processing costs}\par\smallskip
  \small\textbf{Dot plot of sample B's processing costs}\par\smallskip
  \begin{tikzpicture}[x=0.182cm,y=0.29cm,font=\small]
    \draw[black!65] (22,0)--(74,0);
    \foreach \cost in {30,40,50,60,70} {
      \draw[black!65] (\cost,0)--(\cost,-0.3)
        node[below] {\$\cost};
    }
    \foreach \cost/\count in {25/1,34/2,39/3,45/4,48/5,51/4,57/3,62/2,71/1} {
      \foreach \row in {1,...,\count} {
        \fill[taskLinkDarkBlue] (\cost,\row-0.35) circle[radius=1.8pt];
      }
    }
  \end{tikzpicture}
\end{minipage}

There is one transaction at each of \$25 and \$71, two each at \$34 and \$62, three each at \$39 and \$57, four each at \$45 and \$51, and five at \$48. The shape is symmetric with no strong outliers.

We are 95\% confident that the October population mean processing cost lies between \$43.87 and \$52.13 under these assumptions.
\begin{itemize}
  \item \textbf{Design and bias.} Sample randomly from the full October list. Check for missing records and errors in cost or late status.
  \item \textbf{Independence.} Check for dependence such as shared batch costs. Without replacement, check \(n\le 10\%\) of the population. Here \(25\le 1{,}000\).
  \item \textbf{Mean.} With a small sample, look for a roughly symmetric shape without strong outliers, as in sample B. A large sample can support a normal approximation for the mean.
  \item \textbf{Proportion (C5).} Use the one-proportion \(z\) procedure with binary outcomes and at least 10 observations in each category.
\end{itemize}

\Criterion{C8}{Constructs confidence intervals for a population mean or population proportion with unknown variance using the Student’s t-distribution, interpreting the result in finance and economics contexts.}
\Criterion{C9}{Assesses the validity and limitations of confidence-interval inferences by evaluating the sampling design, potential bias, and independence, including the 10\% condition for sampling without replacement and procedure-specific requirements concerning sample size, distributional shape, and outliers.}

\section{Turn an interval into a defensible conclusion}

Using hypothetical sample B, can the business conclude that the population mean processing cost is below its \textbf{\$50 benchmark}?

\VisualHeading{C10 \(\cdot\) From the benchmark question to a conclusion}

{\small\noindent
\begin{tabularx}{\linewidth}{@{}YYYYY@{}}
  \toprule
  \textbf{1. Question} & \textbf{2. Design} & \textbf{3. Procedure}
    & \textbf{4. Estimate and precision} & \textbf{5. Conclusion}\\
  \midrule
  Is \(\mu\) below \$50? & Random sample of 25 & Mean \(t\) interval, \(\sigma\) unknown
    & \(\$48\pm \$4.13\) at 95\% & \$43.87 to \$52.13 spans \$50\\
  \bottomrule
\end{tabularx}\par}

The interval includes values below and above \$50, so the sample leaves the benchmark question unresolved at this confidence level.

The inference covers the business's October online card transactions. In-store transactions use a different processing system and fall outside the sample frame.

\Criterion{C10}{Integrates research design, data analysis, and confidence-interval inference to investigate an economic or financial question, justifying the selected procedure and communicating estimates, precision, assumptions, limitations, and evidence-based conclusions.}

\end{document}
